% Generated by benchmarks/make_bench_tex.py -- do not edit by hand.
% Regenerate with:
%     python benchmarks/crossbench.py
%     python benchmarks/make_bench_tex.py
\documentclass[11pt,a4paper]{article}
\usepackage[utf8]{inputenc}
\usepackage[T1]{fontenc}
\usepackage{lmodern}  % Type1 EC fonts, avoids METAFONT bitmap generation
\usepackage{amsmath,amssymb}
\usepackage[margin=2.7cm]{geometry}
\usepackage{booktabs}
\usepackage{microtype}
\usepackage{xcolor}
\usepackage[colorlinks=true,linkcolor=blue!60!black,citecolor=blue!60!black,urlcolor=blue!60!black]{hyperref}

\title{Benchmarking exact hypervolume computation:\\
       Chan's Section-2 and Section-4.2 algorithms in Python and C}
\author{Michael Emmerich\\[2pt] \normalsize University of Jyv\"askyl\"a, Finland}
\date{\today}

\begin{document}
\maketitle

\begin{abstract}
We compare implementation variants of exact hypervolume computation across six
objective-space dimensions and two families of point sets: the Python
reference implementations of Chan's Section-2 divide and conquer and his
Section-4.2 orthant algorithm, a C99 port of the Section-2 algorithm at two
base-case settings, and three classical baselines --- a dimension sweep after
Guerreiro, Fonseca and Emmerich (CCCG 2012), the WFG algorithm of While,
Bradstreet and Barone (IEEE TEC 2012), and a vendored
Y{\i}ld{\i}z--Suri-style anchored solver. All variants are run on
byte-identical input and agree on the computed hypervolume to within
1.1e-14 relative, which is the primary correctness evidence reported
here. The C port of Chan's algorithm is one to two orders of magnitude faster
than the Python original; among the Python variants, the dimension sweep is
the fastest on every measured case, while the partitioning-based methods'
worst-case guarantees come with visibly larger constants.
\end{abstract}

\section{Setup}

\paragraph{Machine.} Windows-11-10.0.26200-SP0, Python 3.13.15. Times are wall clock, taken in a fresh
subprocess per measurement; fast cases are repeated until the total exceeds
$50$\,ms and the per-call mean is reported. Any measurement exceeding
60.0\,s is recorded as a timeout (\emph{t/o}), and that variant is not
retried at a larger $n$ for the same dimension (the curves are monotone in $n$),
which is shown as \emph{--}.

\paragraph{Reference point.} All objectives are minimised against the reference
point $(1.2,\dots,1.2)$. Datasets are generated from a fixed seed
(20260821), written to disk, and read back by every variant, so all
implementations see exactly the same floating-point input.

\subsection{Variants}

\begin{description}
\item[\textbf{Py \S2}] \texttt{chan\_hypervolume.hypervolume}, the Section-2 divide and conquer, with the $O(n^2d)$ dominance prefilter.
\item[\textbf{Py \S2 np}] The same, with \texttt{prefilter=False}. On these datasets every point is already non-dominated, so the filter can only cost time, never save it.
\item[\textbf{Py \S4.2}] \texttt{chan\_orthant\_dby3.hypervolume\_dby3}, the Section-4.2 orthant algorithm with $\tilde F$ compression on.
\item[\textbf{Py \S4.2 nc}] The same with \texttt{use\_compression=False}, isolating what the compression step buys.
\item[\textbf{C \S2}] The C99 port of Section 2 in \texttt{c/}, base case $b=2$, matching the Python default.
\item[\textbf{C \S2 $b{=}8$}] The same C code with base case $b=8$: the recursion stops earlier and finishes by inclusion--exclusion.
\item[\textbf{Py DS}] \texttt{hv\_baselines.hypervolume\_ds}: dimension sweep with the classical $O(n\log n)$ 3-D base case, organised after Guerreiro, Fonseca, and Emmerich (CCCG 2012) in higher dimensions (exclusive contributions recomputed per point).
\item[\textbf{Py WFG}] \texttt{hv\_baselines.hypervolume\_wfg}: the WFG algorithm of While, Bradstreet, and Barone (IEEE TEC 16(1), 2012), recursive exclusive hypervolumes over dominance-pruned limit sets.
\item[\textbf{Py YS}] The vendored Y{\i}ld{\i}z--Suri-style anchored solver from \texttt{benchmarks/vendor/}, applied through the transform $q = r - y$ (Y{\i}ld{\i}z and Suri, SoCG 2012).
\end{description}

\subsection{Datasets}

\paragraph{\textsf{spherical}.} $n$ points drawn uniformly in the positive
orthant and normalised to the unit sphere. Every point is mutually
non-dominated in all $d$ objectives. This is the standard hard case: dominance
pruning removes nothing, so the full recursion is exercised.

\paragraph{\textsf{cliff}.} The first two coordinates lie on a quarter-circle
arc, and the remaining $d-2$ coordinates are uniform noise. Because the first
two coordinates alone are already mutually non-dominated, no point can dominate
another regardless of the other objectives: the entire dominance structure sits
in a two-dimensional slice, while the remaining $d-2$ objectives contribute
geometry but no order. The front is as large as in the spherical case, but
degenerate in a way the recursion must discover for itself. It is the more
adversarial of the two for any method that hopes to exploit dominance.

\section{Results}

\subsection{Timings}

Table~\ref{tab:spherical} and Table~\ref{tab:cliff} give wall-clock seconds per
call.

\begin{table}[htbp]
\centering
\begin{tabular}{rrrrrrrrrrr}
\toprule
$d$ & $n$ & Py \S2 & Py \S2 np & Py \S4.2 & Py \S4.2 nc & Py DS & Py WFG & Py YS & C \S2 & C \S2 $b{=}8$ \\
\midrule
3 & 100 & 0.0158 & 0.0131 & 0.0319 & 0.0312 & 0.0027 & 0.0195 & 0.0012 & 0.00068 & 0.00077 \\
3 & 200 & 0.0401 & 0.0306 & 0.0761 & 0.0755 & 0.0102 & 0.0725 & 0.0025 & 0.0014 & 0.0015 \\
\addlinespace[2pt]
4 & 60 & 0.0300 & 0.0291 & 0.1661 & 0.1678 & 0.0075 & 0.0314 & 0.0117 & 0.0014 & 0.0011 \\
4 & 120 & 0.0920 & 0.0884 & 0.4070 & 0.4122 & 0.0338 & 0.1079 & 0.0329 & 0.0043 & 0.0041 \\
\addlinespace[2pt]
5 & 40 & 0.0600 & 0.0592 & 0.4023 & 0.4081 & 0.0135 & 0.0509 & 0.1013 & 0.0021 & 0.0011 \\
5 & 80 & 0.2320 & 0.2286 & 1.84 & 1.70 & 0.0607 & 0.3741 & 0.4587 & 0.0071 & 0.0045 \\
\addlinespace[2pt]
7 & 25 & 0.1935 & 0.2056 & 1.35 & 1.33 & 0.0297 & 0.1133 & 1.04 & 0.0087 & 0.0018 \\
7 & 50 & 1.54 & 1.57 & 11.3 & 11.4 & 0.2655 & 2.40 & 11.5 & 0.0510 & 0.0173 \\
\addlinespace[2pt]
9 & 20 & 0.8681 & 0.8689 & 4.93 & 4.76 & 0.0708 & 0.2983 & 17.5 & 0.0315 & 0.0055 \\
9 & 40 & 12.4 & 11.9 & \emph{t/o} & \emph{t/o} & 1.30 & 38.7 & \emph{t/o} & 0.3940 & 0.0730 \\
\addlinespace[2pt]
10 & 20 & 2.31 & 2.37 & 11.1 & 10.8 & 0.1366 & 0.3574 & \emph{t/o} & 0.0750 & 0.0085 \\
10 & 40 & 33.8 & 33.2 & \emph{t/o} & \emph{t/o} & 2.40 & 25.0 & \emph{--} & 1.22 & 0.1990 \\
\addlinespace[2pt]
\bottomrule
\end{tabular}
\caption{\textsf{spherical}: seconds per call.}
\label{tab:spherical}
\end{table}

\begin{table}[htbp]
\centering
\begin{tabular}{rrrrrrrrrrr}
\toprule
$d$ & $n$ & Py \S2 & Py \S2 np & Py \S4.2 & Py \S4.2 nc & Py DS & Py WFG & Py YS & C \S2 & C \S2 $b{=}8$ \\
\midrule
3 & 100 & 0.0135 & 0.0114 & 0.0276 & 0.0293 & 0.0026 & 0.0341 & 0.0014 & 0.00028 & 0.00024 \\
3 & 200 & 0.0354 & 0.0271 & 0.0593 & 0.0598 & 0.0102 & 0.1238 & 0.0031 & 0.0011 & 0.00093 \\
\addlinespace[2pt]
4 & 60 & 0.0162 & 0.0160 & 0.0804 & 0.0823 & 0.0067 & 0.0311 & 0.0090 & 0.00046 & 0.00035 \\
4 & 120 & 0.0518 & 0.0480 & 0.2065 & 0.2163 & 0.0255 & 0.1356 & 0.0225 & 0.0014 & 0.0013 \\
\addlinespace[2pt]
5 & 40 & 0.0307 & 0.0294 & 0.2290 & 0.2025 & 0.0053 & 0.0486 & 0.0828 & 0.0011 & 0.00057 \\
5 & 80 & 0.0915 & 0.0881 & 0.5131 & 0.5112 & 0.0194 & 0.2360 & 0.3169 & 0.0028 & 0.0020 \\
\addlinespace[2pt]
7 & 25 & 0.1213 & 0.1290 & 0.8590 & 0.8570 & 0.0130 & 0.1263 & 1.01 & 0.0048 & 0.0014 \\
7 & 50 & 0.5259 & 0.5288 & 4.11 & 4.17 & 0.0530 & 1.44 & 10.1 & 0.0193 & 0.0056 \\
\addlinespace[2pt]
9 & 20 & 0.4616 & 0.4460 & 3.13 & 3.05 & 0.0220 & 0.0550 & 17.1 & 0.0155 & 0.0023 \\
9 & 40 & 2.78 & 2.76 & 27.6 & 27.9 & 0.1964 & 4.03 & \emph{t/o} & 0.1030 & 0.0168 \\
\addlinespace[2pt]
10 & 20 & 0.9444 & 0.9505 & 5.93 & 5.69 & 0.0356 & 0.1202 & \emph{t/o} & 0.0285 & 0.0036 \\
10 & 40 & 9.39 & 9.34 & \emph{t/o} & \emph{t/o} & 0.3723 & 19.3 & \emph{--} & 0.3210 & 0.0395 \\
\addlinespace[2pt]
\bottomrule
\end{tabular}
\caption{\textsf{cliff}: seconds per call.}
\label{tab:cliff}
\end{table}

\subsection{Speedup of the C port}

Table~\ref{tab:speedup} divides the fastest Python variant for each case by the
C time, so the figures are a lower bound on the speedup over the reference
implementation.

\begin{table}[htbp]
\centering
\begin{tabular}{rrrrrr}
\toprule
 &  & \multicolumn{2}{c}{\textsf{cliff}} & \multicolumn{2}{c}{\textsf{spherical}} \\
\cmidrule(lr){3-4}\cmidrule(lr){5-6}
$d$ & $n$ & $b{=}2$ & $b{=}8$ & $b{=}2$ & $b{=}8$ \\
\midrule
3 & 100 & $\times$5 & $\times$6 & $\times$2 & $\times$2 \\
3 & 200 & $\times$3 & $\times$3 & $\times$2 & $\times$2 \\
\addlinespace[2pt]
4 & 60 & $\times$14 & $\times$19 & $\times$6 & $\times$7 \\
4 & 120 & $\times$16 & $\times$17 & $\times$8 & $\times$8 \\
\addlinespace[2pt]
5 & 40 & $\times$5 & $\times$9 & $\times$6 & $\times$12 \\
5 & 80 & $\times$7 & $\times$10 & $\times$9 & $\times$13 \\
\addlinespace[2pt]
7 & 25 & $\times$3 & $\times$9 & $\times$3 & $\times$17 \\
7 & 50 & $\times$3 & $\times$10 & $\times$5 & $\times$15 \\
\addlinespace[2pt]
9 & 20 & $\times$1 & $\times$9 & $\times$2 & $\times$13 \\
9 & 40 & $\times$2 & $\times$12 & $\times$3 & $\times$18 \\
\addlinespace[2pt]
10 & 20 & $\times$1 & $\times$10 & $\times$2 & $\times$16 \\
10 & 40 & $\times$1 & $\times$9 & $\times$2 & $\times$12 \\
\addlinespace[2pt]
\bottomrule
\end{tabular}
\caption{Speedup of the C port over the fastest Python variant.}
\label{tab:speedup}
\end{table}

\subsection{Recursion size}

Node counts measure the size of the recursion tree itself, so they are
independent of language and of per-node constant factors --- they are the
quantity the two complexity bounds actually bound. Table~\ref{tab:nodes} shows
how the two datasets differ in the work they induce at equal $n$.

\begin{table}[htbp]
\centering
\begin{tabular}{rrrr}
\toprule
$d$ & $n$ & \textsf{cliff} & \textsf{spherical} \\
\midrule
3 & 100 & 417 & 447 \\
3 & 200 & 889 & 943 \\
\addlinespace[2pt]
4 & 60 & 645 & 1\,149 \\
4 & 120 & 1\,779 & 3\,139 \\
\addlinespace[2pt]
5 & 40 & 1\,283 & 2\,463 \\
5 & 80 & 3\,567 & 9\,017 \\
\addlinespace[2pt]
7 & 25 & 4\,647 & 7\,751 \\
7 & 50 & 19\,957 & 57\,093 \\
\addlinespace[2pt]
9 & 20 & 15\,215 & 28\,213 \\
9 & 40 & 95\,481 & 346\,423 \\
\addlinespace[2pt]
10 & 20 & 29\,397 & 65\,613 \\
10 & 40 & 270\,677 & 884\,141 \\
\addlinespace[2pt]
\bottomrule
\end{tabular}
\caption{Recursion nodes, C Section-2 with $b=2$.}
\label{tab:nodes}
\end{table}

Those counts come from the Section-2 recursion in C. The Section-4.2 algorithm
cuts on $(d-3)$-faces rather than $(d-2)$-faces, so its tree should be smaller,
and should grow with an exponent of $d/3$ rather than $d/2$. Running both
solvers over the same instances (\texttt{benchmarks/node\_counts.py}, on a
denser size ladder than the timing runs use) gives
Table~\ref{tab:nodescompare}: the Section-4.2 tree is smaller on every instance
measured, by a factor between 2.0 and 9.1.

\begin{table}[htbp]
\centering
\small
\begin{tabular}{rrrrrrrr}
\toprule
  &   & \multicolumn{3}{c}{\textsf{cliff}} & \multicolumn{3}{c}{\textsf{spherical}} \\
\cmidrule(lr){3-5}\cmidrule(l){6-8}
$d$ & $n$ & $d/2$ & $d/3$ & ratio & $d/2$ & $d/3$ & ratio \\
\midrule
3 & 50 & 175 & 31 & 5.6$\times$ & 229 & 31 & 7.4$\times$ \\
3 & 100 & 417 & 67 & 6.2$\times$ & 455 & 69 & 6.6$\times$ \\
3 & 200 & 891 & 139 & 6.4$\times$ & 933 & 139 & 6.7$\times$ \\
3 & 400 & 2\,345 & 283 & 8.3$\times$ & 1\,947 & 283 & 6.9$\times$ \\
3 & 800 & 5\,193 & 571 & 9.1$\times$ & 3\,891 & 571 & 6.8$\times$ \\
\addlinespace[2pt]
4 & 30 & 245 & 115 & 2.1$\times$ & 427 & 113 & 3.8$\times$ \\
4 & 60 & 651 & 247 & 2.6$\times$ & 1\,141 & 285 & 4.0$\times$ \\
4 & 120 & 1\,771 & 685 & 2.6$\times$ & 3\,143 & 769 & 4.1$\times$ \\
4 & 240 & 5\,279 & 1\,695 & 3.1$\times$ & 7\,109 & 1\,989 & 3.6$\times$ \\
\addlinespace[2pt]
5 & 20 & 431 & 165 & 2.6$\times$ & 573 & 161 & 3.6$\times$ \\
5 & 40 & 1\,283 & 531 & 2.4$\times$ & 2\,477 & 581 & 4.3$\times$ \\
5 & 80 & 3\,541 & 1\,495 & 2.4$\times$ & 9\,017 & 2\,111 & 4.3$\times$ \\
5 & 160 & 12\,549 & 5\,111 & 2.5$\times$ & 26\,945 & 7\,113 & 3.8$\times$ \\
\addlinespace[2pt]
7 & 12 & 559 & 281 & 2.0$\times$ & 1\,015 & 283 & 3.6$\times$ \\
7 & 25 & 4\,647 & 1\,515 & 3.1$\times$ & 7\,751 & 1\,751 & 4.4$\times$ \\
7 & 50 & 19\,959 & 7\,163 & 2.8$\times$ & 57\,093 & 10\,557 & 5.4$\times$ \\
7 & 100 & 75\,673 & 29\,857 & 2.5$\times$ & 316\,667 & 60\,361 & 5.2$\times$ \\
\addlinespace[2pt]
9 & 10 & 933 & 435 & 2.1$\times$ & 1\,775 & 645 & 2.8$\times$ \\
9 & 20 & 15\,215 & 4\,515 & 3.4$\times$ & 28\,195 & 5\,521 & 5.1$\times$ \\
9 & 40 & 95\,453 & 31\,613 & 3.0$\times$ & 346\,421 & 50\,799 & 6.8$\times$ \\
\addlinespace[2pt]
10 & 10 & 1\,727 & 827 & 2.1$\times$ & 2\,765 & 893 & 3.1$\times$ \\
10 & 20 & 29\,397 & 7\,761 & 3.8$\times$ & 65\,613 & 11\,049 & 5.9$\times$ \\
10 & 40 & 270\,681 & 83\,157 & 3.3$\times$ & 884\,199 & 132\,923 & 6.7$\times$ \\
\addlinespace[2pt]
\bottomrule
\end{tabular}
\caption{Recursion nodes, Section 2 versus Section 4.2, same instances. Both
columns are the Python implementations, so they differ from Table~\ref{tab:nodes},
which reports the C port; the two ports resolve ties in the weighted median
slightly differently and so explore marginally different trees.}
\label{tab:nodescompare}
\end{table}

Whether the \emph{exponent} improves is a harder question than whether the tree
is smaller. Table~\ref{tab:nodesexp} fits $\text{nodes} \sim n^{x}$ per
dimension and dataset. Both algorithms stay far below their worst-case
exponents on these instances. The Section-4.2 fit is the lower of the two in
9 of 12 cases, and --- the part that matters --- the gap
widens with $d$, which is where the two bounds diverge: at $d = 10$ the fits are
4.16 against 3.61 on \textsf{spherical}. The exceptions
all sit at low $d$, where $d/2$ and $d/3$ are closest together and the two fits
differ by less than $0.05$ --- within the noise of a fit taken over a four- to
eight-fold range in $n$, at sizes where constant factors still dominate. These
numbers are evidence that the implementations behave as the theory says they
should; they are not a measurement of the asymptotic exponent.

\begin{table}[htbp]
\centering
\small
\begin{tabular}{rrrrrrr}
\toprule
  & \multicolumn{2}{c}{\textsf{cliff}} & \multicolumn{2}{c}{\textsf{spherical}} & \multicolumn{2}{c}{worst case} \\
\cmidrule(lr){2-3}\cmidrule(lr){4-5}\cmidrule(l){6-7}
$d$ & $d/2$ & $d/3$ & $d/2$ & $d/3$ & $d/2$ & $d/3$ \\
\midrule
3 & 1.23 & 1.05 & 1.03 & 1.04 & 1.5 & 1.00 \\
4 & 1.47 & 1.31 & 1.36 & 1.38 & 2.0 & 1.33 \\
5 & 1.61 & 1.64 & 1.85 & 1.83 & 2.5 & 1.67 \\
7 & 2.30 & 2.21 & 2.73 & 2.54 & 3.5 & 2.33 \\
9 & 3.34 & 3.09 & 3.80 & 3.15 & 4.5 & 3.00 \\
10 & 3.65 & 3.33 & 4.16 & 3.61 & 5.0 & 3.33 \\
\bottomrule
\end{tabular}
\caption{Fitted node-count exponents against the worst-case exponents.}
\label{tab:nodesexp}
\end{table}

One further observation from the same runs: the term high-water mark of the
Section-4.2 solver is $1$ on every instance in this table. For grounded
orthants, absorption only adds or merges conditions and never multiplies terms,
so the $\tilde F$ compression --- the only operation that can multiply them ---
never needs to fire.



\subsection{Space}

The recursion is depth first, and each node keeps its own clipped box array
alive while its first child runs, so peak memory is the sum along a
root-to-leaf path rather than the largest single node. Chan's analysis makes
that sum geometric --- the $(d-2)$-face weight falls by $2^{2/d}$ per level ---
and therefore linear. Table~\ref{tab:space} confirms it empirically
(\texttt{benchmarks/space\_profile.py}): the fitted exponent of peak live boxes
in $n$ is 0.88--1.04 across both datasets and
$d = 3, \dots, 7$, with a constant that grows roughly like $2d$ boxes per input
box. A box costs $2d$ doubles, so peak memory is $O(d^2 n)$ words.

\begin{table}[htbp]
\centering
\begin{tabular}{rrrrr}
\toprule
  & \multicolumn{2}{c}{\textsf{cliff}} & \multicolumn{2}{c}{\textsf{spherical}} \\
\cmidrule(lr){2-3}\cmidrule(l){4-5}
$d$ & peak$/n$ & exponent & peak$/n$ & exponent \\
\midrule
3 & 4.2 & 1.00 & 4.4 & 0.98 \\
4 & 5.7 & 0.99 & 6.3 & 0.93 \\
5 & 7.0 & 1.01 & 9.1 & 0.97 \\
6 & 8.5 & 0.89 & 12.1 & 1.04 \\
7 & 10.3 & 0.88 & 15.2 & 1.00 \\
\bottomrule
\end{tabular}
\caption{Peak live boxes per input box, and the fitted exponent in $n$.}
\label{tab:space}
\end{table}

Note that \emph{cumulative} allocation is a different quantity and is
superlinear: every one of the $O(n^{d/2})$ nodes allocates a fresh array. Those
arrays die with their node, so they never coexist.


\section{A family where the Section-4.2 recursion collapses}

The \textsf{wfg\_hard} instances\footnote{From the \texttt{moo-nondominated-sets}
collection, \url{https://github.com/renaudlr/moo-nondominated-sets}.} are built
so that each point is non-dominated through only two of its objectives, the
remaining coordinates being tied. Relative to the root cell every orthant
therefore has exactly two \emph{active} constraints --- and an orthant with at
most two active constraints is precisely what the Section-4.1 simplification
absorbs, as a slab or a 2-sided orthant. The prediction is that the Section-4.2
recursion never cuts at all.

That is exactly what happens (\texttt{benchmarks/wfg\_hard\_probe.py}). On
every instance tested the set of active-constraint counts is exactly
$\{2\}$ --- not merely concentrated there --- so the surviving set $\bar B$ is
empty at the root and the recursion terminates in a single node, for every
dimension and every size:

\begin{table}[htbp]
\centering
\begin{tabular}{rrrrrrr}
\toprule
$d$ & $n$ range & $|\bar B|$ & nodes & absorb (max) & integrate (max) & total exp. \\
\midrule
4 & 100--1000 & 0 & 1 & 0.0026 & 0.0725 & 1.58 \\
6 & 90--900 & 0 & 1 & 0.0038 & 0.0523 & 1.51 \\
8 & 28--300 & 0 & 1 & 0.00095 & 0.0052 & 0.73 \\
10 & 15--150 & 0 & 1 & 0.00055 & 0.0023 & 0.55 \\
\bottomrule
\end{tabular}
\caption{\textsf{wfg\_hard}: the recursion collapses to its base case. Times
are the per-phase maxima over the sizes in the range.}
\label{tab:wfghard}
\end{table}

The whole computation is thus one absorption followed by one integration of a
single basic function, and the phase split shows where the time goes: the
absorption is linear and negligible, while the integration --- Lemma 4.4 ---
dominates and is superlinear, so the total does not reach the $O(n \log n)$
that the collapse would otherwise allow. The engine does not exploit that in
this family the $d/2$ staircase conditions live on \emph{disjoint} axis pairs,
so the integral factorises into $d/2$ independent two-dimensional integrals;
integrating the axes one by one in the general way costs more. This is an
implementation limit rather than an algorithmic one, and it is the clearest
target for future work on this code.

\section{Agreement}

Across 24 cases in which more than one variant completed, the largest
relative spread between the hypervolumes returned by different variants was
1.1e-14, at $d=10$, $n=20$ on \textsf{cliff} (8 variants). Two independent algorithms (Section 2 and
Section 4.2), two languages, and two base-case settings therefore agree to
within floating-point rounding on every case measured. This is the strongest
correctness statement available without an exact-arithmetic oracle, and it is
what the reader should weigh most heavily: the timings below are only
meaningful because all 9 variants compute the same number.

\section{Discussion}

\paragraph{The port is uniformly faster, by a fairly flat factor.} With the
base case matched to the Python default ($b=2$), the C implementation runs
between $\times$1 and $\times$16 faster than the fastest Python
variant on the same input. The factor is notably stable across $d$ and $n$,
which is what one expects when the two programs execute the same recursion and
differ only in interpretation overhead: this is a constant-factor win, not an
algorithmic one.

\paragraph{The base case is the most valuable knob at high dimension.}
Raising the base case to $b=8$ -- stopping the recursion earlier and finishing
by inclusion--exclusion -- widens the gap to between $\times$2 and
$\times$19. The benefit is negligible at $d=3$ and largest at $d=10$.
The recursion pays a per-node cost that grows with $d$ (simplification and cut
selection both sweep all axes), while the base case pays $2^b$ box
intersections regardless of depth; as $d$ grows, trading nodes for a slightly
more expensive base case becomes progressively more favourable. Anyone using
this code at high dimension should tune $b$ rather than accept the default.

\paragraph{Asymptotically faster is not yet actually faster.} The Section-4.2
algorithm has the better bound, $O(n^{d/3}\,\mathrm{polylog}\,n)$, but on every
case measured here it is between $\times$1.7 and $\times$9.9
\emph{slower} than the Section-2 divide and conquer in the same language. Its
constants are large: it manipulates symbolic step functions and terms where
Section~2 manipulates boxes of doubles. The crossover, if it lies within reach
at all, is beyond the problem sizes at which either implementation is usable in
Python. The Section-4.2 code is best read as a faithful and testable
realisation of the paper's construction, not as the method of choice for
computing hypervolumes today. Disabling its compression step changes little at
these sizes, which is consistent with compression being a device for
controlling asymptotic growth rather than a constant-factor optimisation.

\paragraph{The prefilter is free but useless here.} Both datasets consist
entirely of mutually non-dominated points by construction, so the $O(n^2d)$
dominance filter can never remove anything. The measured difference between
\texttt{prefilter=True} and \texttt{prefilter=False} is accordingly small and
of either sign. This says nothing against the filter -- on point sets with
genuine dominance it is a large win -- but it does mean these benchmarks
isolate the recursion rather than the preprocessing.

\paragraph{Degenerate dominance structure is cheaper, not harder.} At equal $n$
and $d$, \textsf{cliff} induces a smaller recursion tree than \textsf{spherical}
(Table~\ref{tab:nodes}) and correspondingly shorter times, even though both
fronts are entirely non-dominated. Concentrating the order structure in two
objectives leaves the remaining $d-2$ coordinates as noise that the
simplification step can often collapse away, whereas a spherical front resists
collapse in every axis at once. The intuition that ``dominance confined to a
low-dimensional slice'' should be a hard case does not survive contact with the
measurements.

\paragraph{Limits of these measurements.} The following exceeded the 60.0\,s limit and are shown as \emph{t/o}: \textsf{cliff} at $d=10$, $n=20$ (Py YS); \textsf{cliff} at $d=10$, $n=40$ (Py \S4.2 nc); \textsf{cliff} at $d=10$, $n=40$ (Py \S4.2); \textsf{cliff} at $d=9$, $n=40$ (Py YS); \textsf{spherical} at $d=10$, $n=20$ (Py YS); \textsf{spherical} at $d=10$, $n=40$ (Py \S4.2 nc); \textsf{spherical} at $d=10$, $n=40$ (Py \S4.2); \textsf{spherical} at $d=9$, $n=40$ (Py YS); \textsf{spherical} at $d=9$, $n=40$ (Py \S4.2 nc); \textsf{spherical} at $d=9$, $n=40$ (Py \S4.2). Timings are
single-run wall clock on one machine with no attempt to control turbo,
thermal state, or background load; the reported factors should be read as
order-of-magnitude, not as precise ratios. Node counts, by contrast, are exact
and machine-independent.

\section{Reproduction}

\begin{verbatim}
make -C c bench_driver
python benchmarks/crossbench.py --timeout 120
python benchmarks/make_bench_tex.py
pdflatex paper/benchmarks.tex
\end{verbatim}

\noindent Add \texttt{-{}-quick} to \texttt{crossbench.py} for the smallest size
per dimension only. All tables and the figures quoted in the discussion are
generated from \texttt{benchmarks/results/crossbench.json}; this document is
produced by \texttt{make\_bench\_tex.py} and should not be edited by hand.

\paragraph{Acknowledgements.}
The benchmark harness and this report were produced by the author, assisted by
Claude Fable~5 (Anthropic), continuing the interactive agentic session that
produced the implementations.

\end{document}
